% Source : cours/cm4.qmd. Risque théorique, pas de causalité empirique ajoutée.
\begin{board}{4}{Loyers : un plafond ne supprime pas la rareté}
 {L'effet réel sur l'offre dépend du dispositif, du marché et de l'horizon.}
\stage{5.4}{1}{Un plafond contraignant}
\stage{16}{2}{Des situations différentes}
\stage{26.6}{3}{L'offre peut s'adapter}
\apartments{4.1}{7.7}
\draw[amber,line width=3pt] (6.4,12)--(9.3,12);
\node[label,text width=3cm] at (7.9,10.7) {Loyer\\plafonné};
\captext{5.4}{Le plafond agit s'il est\\sous le loyer d'équilibre.}
\draw[drawline,fill=blue!7] (11.9,7.3) rectangle (15.2,12);
\person{13.55}{7.7}{blue}
\person{17.4}{7.7}{amber}
\person{19.25}{7.7}{amber}
\captext{16}{Locataire : loyer limité.\\Candidat : accès incertain.}
\noteat{26.6}{12.3}{Court terme}
\apartments{25}{7.2}
\draw[possible] (27.4,10.2)--(29.8,10.2);
\node[smalllabel,text width=3.8cm] at (28.6,7.9) {Autre usage\\possible};
\captext{26.6}{À long terme, un retrait\\de l'offre est possible.}
\end{board}
